% -----------------------------------------------------------------------------
%  Figure 2.1 — Le tableau Scrumban adopté pour la conduite du projet.
%  Le flux continu vient du Kanban (flèche du haut), la revue hebdomadaire vient
%  du Scrum (boucle du bas). La limite de travail en cours est le garde-fou.
% -----------------------------------------------------------------------------
\begin{tikzpicture}[x=1cm, y=1cm]

  % ===================== FLUX CONTINU (Kanban) ===============================
  \draw[fleche, draw=keycemid, line width=1pt] (-7.3,0.9) -- (7.3,0.9);
  \node[etiquette, text=keycemid, font=\sffamily\tiny\bfseries] at (0,0.9)
       {flux continu, hérité du Kanban : une tâche avance dès qu'elle est prête};

  % ===================== LES CINQ COLONNES ===================================
  \foreach \x/\titre/\id in {%
      -6.0/{Backlog}/k1,
      -3.0/{À faire cette semaine}/k2,
       0.0/{En cours}/k3,
       3.0/{En vérification}/k4,
       6.0/{Terminé}/k5}
  {\node[blocfort, text width=2.5cm, anchor=north, minimum height=9mm]
        (\id) at (\x,0.3) {\titre};}

  % --- la limite de travail en cours -----------------------------------------
  \node[bloc, draw=keyceorange, fill=keycesoft, text width=2.5cm, anchor=north,
        minimum height=6mm] (wip) at (0,-0.85)
       {\textbf{limite : 2 tâches}};

  % --- cartes du Backlog -----------------------------------------------------
  \node[bloc, text width=2.5cm, anchor=north] at (-6.0,-0.85)
       {billetterie à code QR};
  \node[bloc, text width=2.5cm, anchor=north] at (-6.0,-1.85)
       {campagnes d'information};
  \node[bloc, text width=2.5cm, anchor=north] at (-6.0,-2.85)
       {carte de partage};

  % --- cartes de la semaine --------------------------------------------------
  \node[bloc, text width=2.5cm, anchor=north] at (-3.0,-0.85)
       {relief 2.5D};
  \node[bloc, text width=2.5cm, anchor=north] at (-3.0,-1.85)
       {généalogie navigable};

  % --- cartes en cours (deux au plus) ----------------------------------------
  \node[bloc, text width=2.5cm, anchor=north] at (0,-1.75)
       {moteur de panorama};
  \node[bloc, text width=2.5cm, anchor=north] at (0,-2.75)
       {assistant conversationnel};

  % --- cartes en vérification ------------------------------------------------
  \node[bloc, text width=2.5cm, anchor=north] at (3.0,-0.85)
       {cloisonnement des données};

  % --- cartes terminées ------------------------------------------------------
  \node[bloc, text width=2.5cm, anchor=north] at (6.0,-0.85)
       {gestion des collections};
  \node[bloc, text width=2.5cm, anchor=north] at (6.0,-1.85)
       {visite immersive 360\textdegree};

  % --- colonnes en fond ------------------------------------------------------
  \begin{scope}[on background layer]
    \foreach \x in {-6.0,-3.0,0.0,3.0,6.0}
      {\draw[rounded corners=4pt, draw=keycegrey!45, dashed, line width=0.5pt]
            (\x-1.42,0.45) rectangle (\x+1.42,-3.65);}
  \end{scope}

  % ===================== CONDITION DE PASSAGE ================================
  \draw[fleche, draw=keyceorange, line width=0.9pt] (4.4,-4.05) -- (7.4,-4.05)
        node[etiquette, pos=0.42, text=keyceorange, anchor=south]
        {passage conditionné par le succès\\ de la chaîne d'intégration continue};

  % ===================== REVUE HEBDOMADAIRE (Scrum) ==========================
  \draw[fleche, draw=keycemid, line width=1pt]
       (7.3,-5.25) -- (-7.3,-5.25);
  \node[etiquette, text=keycemid, font=\sffamily\tiny\bfseries] at (0,-5.25)
       {revue hebdomadaire avec l'encadreur, héritée du Scrum : le reste
        retourne au Backlog};
  \draw[lien, draw=keycemid] (7.3,-3.65) -- (7.3,-5.25);
  \draw[fleche, draw=keycemid] (-7.3,-5.25) -- (-7.3,-3.65);

\end{tikzpicture}
